% Part: computability
% Chapter: computability-theory
% Section: rice-theorem

\documentclass[../../../include/open-logic-section]{subfiles}

\begin{document}

% SEGMENT OLP-0248-S01

\olfileid{cmp}{thy}{rce}
\olsection{ரைஸின் தேற்றம்}

சிந்தித்துப் பார்த்தால், \olref[tot]{prop:total} இன் நிரூபணத்தில்
$\fn{Tot}$ இன் குறிப்பான அம்சங்கள் பங்கு வகிக்கவில்லை என்பது தெரியும்.
$x$ ஆனது $K$ இல் இருக்கும்போது சரியாக மாறிலிச் சார்பு $j(y) = 0$
போல நடந்துகொள்ளுமாறு $h(x,y)$ ஐ வடிவமைத்தோம்; ஆனால் அச்சூழலில்
வேறு எந்தப் பகுதி கணிக்கத்தக்க சார்புபோலவும் அது நடந்துகொள்ளுமாறு
அமைத்திருக்கலாம். கணிக்கத்தக்க சார்புகளின் அற்பமல்லாத எந்தப்
பண்பும் தீர்மானிக்கத்தக்கது அல்ல என்று தோராயமாகக் கூறும் மேலும்
பொதுவான தேற்றத்தை முன்வைக்க இக்கவனிப்பு உதவுகிறது.

$\cfind{0}$, $\cfind{1}$, $\cfind{2}$,~\dots என்பது பகுதி
கணிக்கத்தக்க சார்புகளின் நமது செந்தர எண்ணீடு என்பதை நினைவில் கொள்க.

\begin{thm}[ரைஸின் தேற்றம்]
  $C$ பகுதி கணிக்கத்தக்க சார்புகளின் ஏதாவது ஒரு கணமாகவும்,
  $A = \Setabs{n}{\cfind{n} \in C}$ ஆகவும் இருக்கட்டும்.
  $A$ கணிக்கத்தக்கது எனில், $C$ வெற்றுக்கணம் அல்லது $C$ ஆனது பகுதி
  கணிக்கத்தக்க சார்புகள் அனைத்தின் கணம்.
\end{thm}

{\em சுட்டெண் கணம்} என்பது பின்வரும் பண்புள்ள கணம் $A$: ஒரே சார்பைக்
``கணிக்கும்'' சுட்டெண்களாக $n$, $m$ இருந்தால், $n$, $m$ இரண்டுமே
$A$ இல் இருக்கும் அல்லது இரண்டுமே இருக்காது. தேற்றத்திலுள்ள கணம்~$A$
க்கு இப்பண்பு இருப்பதைக் காண்பது கடினமல்ல. மறுதிசையில், $A$ ஒரு
சுட்டெண் கணமாகவும், இச்சுட்டெண்களால் கணிக்கப்படும் சார்புகளின் கணமாக
~$C$ உம் இருந்தால், $A = \Setabs{n}{\cfind{n} \in C}$.

\begin{explain}
இச்சொல்லாட்சியில், அற்பமல்லாத எந்தச் சுட்டெண் கணமும்
தீர்மானிக்கத்தக்கது அல்ல என்ற கூற்றுக்குச் சமானமானதே ரைஸின் தேற்றம்.
தேற்றம் கூறுவதைப் புரிந்துகொள்ள, \emph{நிரல்கள்} (உங்களுக்கு விருப்பமான
நிரலாக்க மொழியில் உள்ளவை என்க) மற்றும் அவை கணிக்கும் சார்புகள்
ஆகியவற்றுக்கு இடையிலான வேறுபாட்டை வலியுறுத்துவது உதவும்.
தொடரியல் பொருட்களான நிரல்களைப் (சுட்டெண்களைப்) பற்றிக் கணிக்கத்தக்க
கேள்விகள் நிச்சயமாக உள்ளன: இந்நிரலில் 150-க்கும் மேற்பட்ட குறியீடுகள்
உள்ளனவா? 22-க்கும் மேற்பட்ட வரிகள் உள்ளனவா? ``while'' கூற்று
உள்ளதா? ``print'' கூற்றுக்கு அளிக்கப்படும் மதிப்புருவாக
``hello world'' என்ற எழுத்துத்தொடர் எப்போதாவது தோன்றுகிறதா?
நிரலின் \emph{நடத்தை} பற்றிய அற்பமல்லாத எந்தக் கேள்வியும்
கணிக்கத்தக்கது அல்ல என்று ரைஸின் தேற்றம் கூறுகிறது. பின்வருவனவும்
அத்தகைய கேள்விகளே: உள்ளீடு $0$ இல் நிரல் நிற்கிறதா? அது
எப்போதாவது நிற்கிறதா? அது எப்போதாவது ஓர் இரட்டை எண்ணை வெளியிடுகிறதா?
\end{explain}

\begin{proof}[ரைஸின் தேற்றத்தின் நிரூபணம்]
$C$ வெற்றுக்கணமும் அல்ல, பகுதி கணிக்கத்தக்க சார்புகள் அனைத்தின்
கணமும் அல்ல எனக் கொள்க; மேலும்~$C$ இலுள்ள சார்புகளின் சுட்டெண்களின்
கணமாக~$A$ இருக்கட்டும். $A$ கணிக்கத்தக்கதாக இருந்தால் நிற்றல்
சிக்கலைத் தீர்க்க முடியும் என்று காட்டுவோம்; எனவே $A$ கணிக்கத்தக்கது அல்ல.

பொதுத்தன்மை இழப்பின்றி, எங்குமே வரையறுக்கப்படாத சார்பு $f$ ஆனது
$C$ இல் இல்லை எனக் கொள்ளலாம் (அவ்வாறு இல்லாவிட்டால், கீழுள்ள
வாதத்தில் $C$ ஐயும் அதன் நிரப்பியையும் இடமாற்றுக). $g$ ஆனது~$C$
இலுள்ள ஏதாவது ஒரு சார்பு ஆகட்டும். $A$ ஐத் தீர்மானிக்க முடிந்தால்,
$f$ ஐக் கணிக்கும் சுட்டெண்களையும் $g$ ஐக் கணிக்கும் சுட்டெண்களையும்
வேறுபடுத்திக் கூறலாம்; பிறகு அத்திறனைப் பயன்படுத்தி நிற்றல் சிக்கலைத்
தீர்க்கலாம் என்பதே கருத்து.

முறை இதுவே. அனைத்தளாவிய பகுதி கணிக்கத்தக்க சார்புகளைப் பயன்படுத்தி,
ஒரு சார்பைப் பின்வருமாறு வரையறுக்கலாம்:
\[
h(x,y) \simeq
\begin{cases}
\text{வரையறுக்கப்படவில்லை} & \text{$\cfind{x}(x) \fundefined$ எனில்} \\
g(y) & \text{இல்லையெனில்.}
\end{cases}
\]
$h$ ஐக் கணிக்க, முதலில் $\cfind{x}(x)$ ஐக் கணிக்க முயல்கிறோம்;
அக்கணிப்பு நின்றால் தொடர்ந்து~$g(y)$ ஐக் கணிக்கிறோம்; மேலும்
{\em அந்தக்} கணிப்பும் நின்றால் வெளியீட்டைத் திருப்பித்தருகிறோம்.
மேலும் முறைசார்ந்து,
\[
h(x,y) \simeq \Proj{2}{0}(g(y),\fn{Un}(x,x)).
\]
என்று எழுதலாம். இங்கு $\Proj{2}{0}(z_0, z_1) = z_0$ என்பது
$0$-ஆம் மதிப்புருவைத் திருப்பித்தரும் கணிக்கத்தக்க $2$-இட வீழ்த்தல் சார்பு.

அப்போது $h$ பகுதி கணிக்கத்தக்க சார்புகளின் ஒரு சேர்ப்பு. வலப்பக்கம்
~$\fn{Un}(x,x)$ உம் $g(y)$ உம் வரையறுக்கப்பட்டிருக்கும்போது மட்டுமே
வரையறுக்கப்பட்டு~$g(y)$ குச் சமமாக உள்ளது.

ஒரு நிலையான~$x$ க்கு, $\cfind{x}(x)$ வரையறுக்கப்படவில்லை எனில்
ஒவ்வொரு~$y$ க்கும் $h(x,y)$ வரையறுக்கப்படவில்லை; $\cfind{x}(x)$
வரையறுக்கப்பட்டிருந்தால் $h(x,y) \simeq g(y)$ என்பதைக் கவனிக்க.
எனவே எந்த நிலையான~$x$ மதிப்புக்கும், $h(x,y)$ ஆனது $f$ போலவே
அல்லது $g$ போலவே நடந்துகொள்கிறது; $\cfind{x}(x)$ வரையறுக்கப்பட்டுள்ளதா
இல்லையா என்று தீர்மானிப்பது, இவ்விரு நேர்வுகளில் எது நிலவுகிறது
என்று தீர்மானிப்பதற்குச் சமம். ஆனால் இது
$h_x(y) \simeq h(x,y)$ ஆனது~$C$ இல் உள்ளதா இல்லையா என்று
தீர்மானிப்பதற்குச் சமம்; $A$ கணிக்கத்தக்கதாக இருந்தால் அதையே
நம்மால் செய்ய முடியும்.

மேலும் முறைசார்ந்து, $h$ பகுதி கணிக்கத்தக்கதாக இருப்பதால்,
ஏதோ சுட்டெண்~$e$ க்கு அது சார்பு $\cfind{e}$ குச் சமம்.
$s$-$m$-$n$ தேற்றத்தால், ஒவ்வொரு~$x$ க்கும்
$\cfind{s(e,x)}(y) = h_x(y)$ என அமையும் முதன்மை மீள்வரையறைச்
சார்பு~$s$ ஒன்று உள்ளது. இப்போது ஒவ்வொரு $x$ க்கும்,
$\cfind{x}(x) \fdefined$ எனில் $\cfind{s(e,x)}$ ஆனது~$g$ என்ற
அதே சார்பு; எனவே $s(e,x)$ ஆனது $A$ இல் உள்ளது. மறுபுறம்,
$\cfind{x}(x) \uparrow$ எனில் $\cfind{s(e,x)}$ ஆனது~$f$ என்ற
அதே சார்பு; எனவே $s(e,x)$ ஆனது~$A$ இல் இல்லை. வேறுவிதமாகக்
கூறினால், ஒவ்வொரு~$x$ க்கும், $x \in K$ எனில், எனில் மட்டுமே,
$s(e,x) \in A$. $A$ கணிக்கத்தக்கதாக இருந்தால் $K$ உம் அவ்வாறே
இருக்கும்; இது முரண்பாடு. எனவே $A$ கணிக்கத்தக்கது அல்ல.
\end{proof}

ரைஸின் தேற்றம் மிகவும் வலிமையானது. பின்வரும் உடனடி முடிவுரை சில
மாதிரிப் பயன்பாடுகளைக் காட்டுகிறது.
\begin{cor}
பின்வரும் கணங்கள் தீர்மானிக்கவியலாதவை.
\begin{enumerate}
\item $\Setabs{x}{\text{$17$ ஆனது $\cfind{x}$ இன் வீச்சில் உள்ளது}}$
\item $\Setabs{x}{\text{$\cfind{x}$ மாறிலியாகும்}}$
\item $\Setabs{x}{\text{$\cfind{x}$ முழுமையானது}}$
\item $\Setabs{x}{\text{$y < y'$, $\cfind{x}(y) \fdefined$ ஆகவும்,
    $\cfind{x}(y') \fdefined$ எனில் $\cfind{x}(y) < \cfind{x}(y')$
    ஆகவும் இருக்கும் போதெல்லாம்}}$
\end{enumerate}
\end{cor}

\begin{proof}
  இவை அனைத்தும் அற்பமல்லாத சுட்டெண் கணங்கள்.
\end{proof}

\end{document}
